\section{Conclusion}

Closed-set acoustic threat classifiers do not survive a change of forest: strong
in-domain validation hides a cross-site collapse once a deployable
false-positive budget is enforced, and the collapse is one of operating-point
transfer rather than discrimination, driven by a representation that encodes
recording site more strongly than event class. We reframed the task as
background-first and open-set, flagging deviation from a site's own soundscape,
attributing with incrementally built prototypes, and abstaining honestly. The
resulting detector ships with zero labeled threats and improves as rangers
confirm alerts, without retraining.

Across four held-out sites the choice of frozen representation, not the detector
head, governed transfer, and a supervised probe on the same embeddings priced
what label-free operation costs: labels roughly double recall, but only the
label-free detector can run on day one. BirdNET sharpened the central lesson by
being the strongest probe substrate while flagging nothing under a transferred
threshold, a reminder that ranking quality and deployability are different
properties and that false-positive control is what binds.

Beyond forests, the same recipe fits any monitoring setting where normal is easy
to record, anomalies are rare and open-ended, and expert attention is the
limiting resource. Limitations remain: one threat class, four sites, two of them
small, modest absolute recall, and a calibration gap that a little
destination-site background should close. Field validation on in-situ positives
is the necessary next step and is underway.
